\section{Methodology}
\label{sec:method}

\subsection{AR Plan Generation}

The \ar{} plan for each state is generated by the \texttt{bisect} binary
at the minimum-edge-cut seed found by \cs{} at threshold $T = 600$
\citep{dellaLibera2026b16}.  Full details of the algorithm are in
\citet{dellaLibera2026b11}.  Briefly: the seat count $k$ is factored into
primes; each prime factor $p$ generates a $p$-way minimum-cut partition of
the current subgraph; the hierarchy continues until every leaf contains
exactly one district.  Where $k$ is prime (Pennsylvania, $k = 17$), a
single $k$-way partition is used.

\subsection{Ensemble Generation via GerryChain ReCom}

Ensembles are generated using \gc{} \recom{} \citep{deford2021recombination}.
\recom{} produces plans uniformly at random from the feasible space
(population-balanced, contiguous plans).  Each run uses 1,000 steps.
At N=1,000 steps, the ensemble provides a preliminary distributional portrait;
multi-seed validation (Phase 2) is required to confirm mixing.
The ensemble provides an empirical distribution of the normalised edge-cut
fraction $\EC$ for each state.

\begin{definition}[Normalised edge-cut fraction]
For a redistricting plan $\pi$ on census-tract graph $G = (V, E)$:
\begin{equation}
  \EC(\pi) = \frac{|\{(u,v) \in E : \pi(u) \neq \pi(v)\}|}{|E|}.
\end{equation}
This is the fraction of tract-adjacency edges that cross district boundaries.
Lower $\EC$ corresponds to higher compactness: districts share fewer boundaries
with each other.
\end{definition}

\subsection{Percentile Placement}

The \ar{} plan's percentile on $\EC$ is computed as:
\begin{equation}
  \hat\pi_{\EC} = \frac{|\{i : \EC_i < \EC_{\mathrm{AR}}\}|}{N},
\end{equation}
where $\EC_1, \ldots, \EC_N$ are the ensemble plan edge-cut fractions.
A \emph{low} $\hat\pi_{\EC}$ means the \ar{} plan has lower cut than most
ensemble plans---i.e., it is more compact.

\subsection{Why Edge Cut, Not Polsby-Popper}

Polsby-Popper (PP) measures the ratio of district area to perimeter squared.
It is a geometric property of the final district shapes.  Edge cut measures
the number of shared boundary segments between districts.  METIS minimises
edge cut directly; the resulting plans tend to have high PP but the
relationship is not one-to-one.  We report edge-cut percentiles because:
\begin{enumerate}
  \item Edge cut is the \ar{} algorithm's direct objective;
  \item The \recom{} ensemble provides edge-cut statistics for direct comparison;
  \item Courts can audit edge cut from the assignment file alone (no geometry
    computation required), supporting verifiability.
\end{enumerate}

\subsection{Uncertainty Quantification}

\paragraph{Sampling uncertainty for G.1's 1{,}000-step runs.}
Sampling uncertainty is substantial at $N = 1{,}000$ steps.
Using the lag-1 autocorrelation estimates from G.4~\citep{dellaLibera2026g4}
($\rho_1(\mathrm{PP}) \approx 0.87$ for NC-scale chains), the effective sample
size for G.1's ensembles is approximately:
\[
  \ess \approx \frac{N(1-\rho_1)}{1+\rho_1}
  = \frac{1{,}000 \times 0.13}{1.87} \approx 70.
\]
The standard error on the 0.2nd-percentile compactness claim is:
\begin{equation}
  \mathrm{SE}(\mathrm{PP}) = \sqrt{\hat\pi(1-\hat\pi)/\ess}
  = \sqrt{0.002 \times 0.998 / 70} \approx 0.0053,
\end{equation}
or about 0.5 percentage points for the compactness point estimate.
For Democratic seat-count percentiles (D-seats $\rho_1 \approx 0.75$,
$\ess_D \approx 143$), the 90\% confidence interval at the median is
approximately $\pm\,6.9$ percentage points.
\emph{All G.1 percentile figures are preliminary point estimates;
multi-chain Phase~2 validation is required before these estimates
can be cited as certified ensemble positions.}

\paragraph{Herschlag chain (reference only).}
For context, the NC Herschlag chain ($N = 24{,}518$) yields
$\ess \approx 1{,}703$, giving a 90\% CI of approximately
$\pm 2.0$ percentile points at the seat-count median.
This ESS applies to the Herschlag chain and should not be used to
certify G.1's 1{,}000-step results.

For Texas and California (pending Phase~2), we report ``$< 5$th percentile''
as a design-based lower bound pending full ensemble runs.
